\documentclass[10pt,a4paper]{amsart}
\usepackage[margin=19mm]{geometry}
\usepackage{amsmath,amssymb,mathtools}
\usepackage[scr=rsfs,cal=euler]{mathalfa}
\usepackage{xfrac}
\usepackage[unicode,hidelinks,bookmarks=false]{hyperref}
\newcommand{\ie}{i.e.\ }
\newcommand{\cf}{cf.\ }
\newcommand{\resp}{respectively\ }
\newcommand{\Subsection}{\subsection}
\providecommand{\nobreakdash}{\nobreak\hbox{-}}
\setlength{\emergencystretch}{2em}
\multlinegap0pt
\usepackage{bm}
\usepackage{bbm}
\usepackage{dsfont}
\usepackage{xspace}
\usepackage{cancel}
\usepackage{mathtools}
\usepackage{cleveref}
\usepackage{amsthm}
\makeatletter
\@ifundefined{proposition}{\newtheorem{proposition}{Proposition}}{}
\makeatother
\DeclareMathOperator{\conv}{conv}
\newcommand\KG[2]{K_G\left({#1}\ifx#2\empty\else\mapsto\fi{#2}\right)}
\newcommand\cut[2]{\mathrm{SDP}_{#1}^{#2}}
\newcommand\sdp[2]{\mathrm{SDP}_{#1}\ifx#2\empty\else\left({#2}\right)\fi}
\newcommand{\innp}[2]{\left\langle{#1},{#2}\right\rangle}
\newcommand{\ta}{\tilde{a}}
\newcommand{\tb}{\tilde{b}}
\newcommand{\tinnp}[2]{\langle{#1},{#2}\rangle}
\renewcommand{\geq}{\geqslant}
\renewcommand{\leq}{\leqslant}
\title[Better bounds on finite-order Grothendieck constants]{Better bounds on finite-order Grothendieck constants}
\author{S\'ebastien Designolle, Tam\'as V\'ertesi, Sebastian Pokutta}
\date{Source: arXiv:2409.03739v3 (CC BY 4.0)}
\begin{document}
\maketitle

\noindent\emph{Source and licence.} Better bounds on finite-order Grothendieck constants. Version arXiv:2409.03739v3 (CC BY 4.0), distributed under the Creative Commons Attribution 4.0 International licence, as recorded in the arXiv licence metadata for this exact version and shown on the abstract page https://arxiv.org/abs/2409.03739v3. Full rights notice: licenses/arxiv-2409.03739-CC-BY-4.0.txt.\par\smallskip

\noindent\textit{Editorial scope.} Counted unit: the Proposition whose source label is \texttt{eqn:upper} in the article (source0.tex lines 561--573, source label \texttt{eqn:upper}), delivered together with the article's own complete original proof of it (source0.tex lines 575--603, one \texttt{proof} environment of 29 lines). No mathematics is written, repaired or re-derived: statement and proof are the article's own text line for line. Required context retained uncounted: eqn:SDPdM (lines 122--125). Every definition, display and label of those lines is retained unchanged. Omitted from this package and not redistributed: 1--121, 126--560, 574--574, 604--683. Nothing inside a retained excerpt is omitted or altered: every retained body is delivered exactly as the article's lines are written, so no omission receipt of any kind accompanies this package. Citations: the retained excerpts carry no citation command at all, so no bibliography entry of the article is transcribed and no reference is redistributed. Reference adaptation: none was needed and none was made; every reference inside the delivered unit resolves inside it, so no unresolved reference or citation is produced. Printed slips kept literal: no printed word, symbol, hypothesis or number of the counted statement or of its proof was altered. Layout and class: the article's own class is \texttt{revtex4-1} with packages including algorithm, algpseudocode, amsfonts, amsmath, amssymb, amsthm, colortbl, dsfont, graphicx, hyperref, hyphenat, lmodern, mathtools, microtype; the wrapper is the lane's pinned \texttt{amsart} editorial wrapper at 10pt on A4, which restates the theorem-like environments the retained text needs and adds the article's own macro definitions that the retained text needs (\texttt{\textbackslash KG}, \texttt{\textbackslash conv}, \texttt{\textbackslash cut}, \texttt{\textbackslash geq}, \texttt{\textbackslash innp}, \texttt{\textbackslash leq}, \texttt{\textbackslash sdp}, \texttt{\textbackslash ta}, \texttt{\textbackslash tb}, \texttt{\textbackslash tinnp}), with extra wrapper packages (bm, bbm, dsfont, xspace, cancel, mathtools, cleveref). The delivered numbering is the unit's own. Assets: no retained span contains a figure environment, a graphics-inclusion command, a PDF-page inclusion, a table or a TikZ picture; no image file is copied into this package, the package declares no asset dependency and no image file is redistributed with it. Licence: version arXiv:2409.03739v3 of this article is distributed under the Creative Commons Attribution 4.0 International licence, as recorded in the arXiv licence metadata for this exact version and shown on the abstract page https://arxiv.org/abs/2409.03739v3. A full rights notice is delivered at licenses/arxiv-2409.03739-CC-BY-4.0.txt. Characters: every retained excerpt is pure ASCII, so the delivered engine, XeLaTeX with its default Latin Modern text font, reports no missing character.
\begin{equation}
  \sdp{d}{M}=\max\left\{\sum_{x=1}^{m_1}\sum_{y=1}^{m_2} M_{xy}\innp{a_x}{b_y}~\Big|~\forall x\in[m_1],~a_x\in S^{d-1},~\forall y\in[m_2],~b_y\in S^{d-1}\right\},
  \label{eqn:SDPdM}
\end{equation}

\begin{proposition}
  Let $P\in\cut{d}{p_1,p_2}$ with underlying vectors $a_1^P,\ldots,a_{p_1}^P,b_1^P,\ldots,b_{p_2}^P\in S^{d-1}$ such that
  \begin{itemize}
    \item there exist $\eta_1,\eta_2\in(0,1)$ for which $\eta_1S^{d-1}\subset\conv\{a_i^P\}_{i=1}^{p_1}$ and $\eta_2S^{d-1}\subset\conv\{b_j^P\}_{j=1}^{p_2}$,
    \item there exists $\alpha\in(0,1)$ such that $\alpha P\in\cut{n}{p_1,p_2}$.
  \end{itemize}
  Then we have the following bound:
  \begin{equation}
    \KG{d}{n}\leq\frac{1}{\alpha\,\eta_1\eta_2}.
    \label{eqn:upper}
  \end{equation}
  \label{prop:upper}
\end{proposition}

\begin{proof}
  Consider a matrix $M$ of size $m_1\times m_2$.
  Since
  \begin{equation}
    \alpha\,\eta_1\eta_2\,\sdp{d}{M}=\alpha\,\eta_1\eta_2\max_{N\in\cut{d}{m_1,m_2}}\innp{M}{N}=\max_{N\in\cut{d}{m_1,m_2}}\innp{M}{\alpha\,\eta_1\eta_2\,N},
    \label{eqn:proof_upper1}
  \end{equation}
  we focus on $\alpha\,\eta_1\eta_2\,N$ for $N\in\cut{d}{m_1,m_2}$ with underlying vectors $a_1^N,\ldots,a_{m_1}^N,b_1^N,\ldots,b_{m_2}^N\in S^{d-1}$.

  By the first assumption on $\{a_i^P\}_{i=1}^{p_1}$ and $\{b_j^P\}_{j=1}^{p_2}$, we know that we can write, for $x\in[m_1]$ and $y\in[m_2]$,
  \begin{equation}
    \eta_1a_x^N=\sum_{i=1}^{p_1}v_i^{(x)}a_i^P\quad\text{and}\quad\eta_2b_y^N=\sum_{j=1}^{p_2}w_j^{(y)}b_j^P,\quad\text{so that}\quad\alpha\,\eta_1\eta_2\,N_{xy}=\alpha\sum_{i,j}v_i^{(x)}w_j^{(y)}\innp{a_i^P}{b_j^P}
  \end{equation}
  where $v_i^{(x)}\geq0$, $w_j^{(y)}\geq0$, $\sum_iv_i^{(x)}=1$, and $\sum_jw_j^{(y)}=1$.

  Next we use the second assumption, which can be written as follows: there exist positive weights $\lambda_k$ summing up to one and vectors $\ta_1^{(k)},\ldots,\ta_{p_1}^{(k)},\tb_1^{(k)},\ldots,\tb_{p_2}^{(k)}\in S^{n-1}$ indexed by $k$ (in a finite set) such that $\alpha\tinnp{a_i^P}{b_j^P}=\sum_k\lambda_k\tinnp{\ta_i^{(k)}}{\tb_j^{(k)}}$ holds for all $i\in[p_1]$ and $j\in[p_2]$.
  This gives, for $x\in[m_1]$ and $y\in[m_2]$,
  \begin{equation}
    \alpha\,\eta_1\eta_2\,N_{xy}=\sum_{i,j}v_i^{(x)}w_j^{(y)}\sum_k\lambda_k\innp{\ta_i^{(k)}}{\tb_j^{(k)}}=\sum_k\lambda_k\innp{\sum_iv_i^{(x)}\ta_i^{(k)}}{\sum_jw_j^{(y)}\ta_j^{(k)}}=\sum_k\lambda_k\innp{\ta_x^{(k)}}{\tb_y^{(k)}},
    \label{eqn:proof_upper2}
  \end{equation}
  where we have defined the vectors $\ta_x^{(k)}=\sum_iv_i^{(x)}\ta_i^{(k)}$ and $\tb_y^{(k)}=\sum_jw_j^{(y)}\ta_j^{(k)}$ in the $n$-dimensional unit ball $B^n$.

  Combining \cref{eqn:proof_upper1} with \cref{eqn:proof_upper2} and using the definition of $\sdp{d}{M}$ in \cref{eqn:SDPdM}, we obtain
  \begin{equation}
    \alpha\,\eta_1\eta_2\,\sdp{d}{M}\leq\max\left\{\sum_{x=1}^{m_1}\sum_{y=1}^{m_2} M_{xy}\innp{\ta_x}{\tb_y}~\Big|~\forall x\in[m_1],~\ta_x\in B^n,~\forall y\in[m_2],~\tb_y\in B^n\right\}\leq\sdp{n}{M},
  \end{equation}
  so that \cref{eqn:upper} follows from the definition of $\KG{d}{n}$.
\end{proof}

\end{document}
